\section{Conclusion}
\label{sec:conclusion}

This paper has presented the design, implementation, and qualitative
evaluation of a cabinet-scale prototype for household food-waste
reduction. The prototype integrates four HX711-amplified load cells
for shelf-weight estimation, a DHT11 for temperature and humidity, an
MQ135 for indicative gas observation, an AI Thinker ESP32-CAM for
QR-based product registration, and an LED-strip illumination
subsystem within a single low-cost integrated platform. Sensor output
is transported from an Arduino Uno over USB to a Python bridge that
exposes the readings as a CORS-enabled local HTTP endpoint, and a
cross-platform Flutter application provides dashboard, camera,
product, alert, and configuration views over the same feed.

The principal contributions are the multi-modal integration of four
sensing channels in one reference design; the USB-to-HTTP bridge,
which decouples the sensor microcontroller from the wireless network;
the tiled multi-QR routine, which recovers several stickers from a
single capture; the pixel-level RGB ripeness classifier, which serves
as an interpretable baseline against which model-based detectors can
be compared; and the public archival of the complete implementation
in the Software Heritage Archive, supporting independent
reproduction.

Demonstration sessions confirmed concurrent operation of all
subsystems with end-to-end latency below one second, reliable weight
detection above approximately 50\,g, and reproducible visual ripeness
estimation. The limitations documented in
Section~\ref{sec:discussion} map directly to a programme of future
work: hardware improvements for lightweight-product sensitivity, a
learned visual classifier for foods beyond bananas, a multi-shelf
architecture, and a server-mediated synchronisation layer suitable
for deployments that cross network boundaries.
